\documentclass[11pt]{article}
\usepackage{climbingguide}

\begin{document}
\CGSetHeaderLeft{Serra do Cuó}

\CGCoverPage{book_cover_front.jpg}{SERRA DO CUÓ}{{\CGCoverSubtitleFont Total de Vias:} \CGTotalRoutes}

\CGTextSingleColumn{}{
    Guia de Escalada da Serra do Cuó\\
    %Autores: N/A\\
    %Revisão: N/A
    Fotos: Acervo Brito, CuóRockTrip2026, Thibor, Lu. \\
    Editor: Lu Freitas\\
    Edição: \CGToday \\
    Número de Vias: \CGTotalRoutes \\
    Agradecimentos: Prefeitura de Campo Grande RN, AERN, EENe. \\ Agradecimento especial à Agnaldo por sempre apoiar, bem como à João Victor e a todos que se engajaram em fomentar à escalada em Serra do Cuó.
}
\vspace{5mm}
\begin{center}
    {\huge EDIÇÃO NÃO OFICIAL \\[5pt] EM DESENVOLVIMENTO }
\end{center}
\vfill
\CGZoneStats{Global}
\CGTextSingleColumn{}{
    Número de vias por graduação.
}

\CGSetHeaderLeft{SERRA DO CUÓ}
\CGSetHeaderRight{Índice}
\CGTableOfContents
\clearpage


\CGSetHeaderRight{APRESENTAÇÃO}

%\CGTextSingleColumn{Apresentação}{
%Inspirado em grandes encontros da comunidade e no constante desenvolvimento de novos picos de rocha no Brasil, este guia nasce do esforço conjunto para consolidar a escalada na Serra do Cuó, em Campo Grande-RN.

%Mais do que um catálogo, este material tem como objetivo auxiliar a comunidade no reconhecimento do local, organizar o acesso aos diversos setores e, acima de tudo, promover a integração sustentável entre os escaladores e a população local. Todo o desenvolvimento esportivo na região contou, ao longo do tempo, com o fundamental apoio da Prefeitura Municipal de Campo Grande e da Associação de Escaladores do Rio Grande do Norte (AERN).

%Preparamos este guia para que ele seja a base definitiva das suas aventuras na Serra do Cuó. Que todos tenham excelentes escaladas!}

%\CGDoubleTextSection{Primeira Impressão: O Despertar para o Cuó}{
%\textit{Por Flávia dos Anjos}

%Minha primeira impressão da Serra do Cuó foi incrível. A vista de longe é de duas enormes paredes laranja que só me lembram as imagens que vi das tão aclamadas paredes da Catalunha espanhola. Já passavam das 9h30 da manhã e o sol não estava de brincadeira. A trilha era claramente longa e exposta, o que me deixou com um sentimento misto de “eu preciso ir nessa linda parede AGORA!” intercalado com “Eu não quero subir essa trilha de jeito NENHUM!”.

%Para minha sorte (ou não), os amigos Janine Falcão e Allysson Laurentino tinham programado algo bem mais amigável para o pouco tempo que tínhamos disponível. Acontece que, ao longo da inclinada vegetação que se forma aos pés dos dois paredões, a superfície é recoberta de enormes blocos de pedra. Muitos blocos, com todas as faces verticais e negativas, oferecendo inúmeras vias e permitindo que os escaladores escolhessem o lado da sombra e se ajustassem a ela ao longo do dia.

%O que mais me chamou a atenção foi a existência de fendas, muitas fendas. Escalei algumas vias ali e no fim me deparei com uma linha que me chamou muita atenção. Perguntei que via era e o Allysson me confirmou que era um projeto; o conquistador tinha montado a linha, mas não tinha conseguido a tão desejada “cadena”.
%}{}{
%Me interessei e resolvi tentar. O grau da via concentra-se nos dois movimentos iniciais, um \textit{boulder} de calcanhar seguido de um movimento dinâmico. A cadena só saiu depois de um descanso, um lanchinho e um terceiro e último “pega”. Virando o \textit{boulder} a via fica bem amigável, a ser protegida com peças móveis médias.

%Já que o Allysson me garantiu que eu estava levando a “primeira ascensão” (FA), ele me deu o direito de batizar a via, e eu sugeri “Ainda não mijei hoje”.

%Depois de um dia no Cuó eu olhei para trás de novo e voltei a encarar os dois paredões. Refleti sobre o local, sobre aquele projeto sem cadena e terminei com minha segunda e última impressão do Cuó: muito potencial! Uma rocha magnífica, cheia de fendas, com oportunidades para todos os níveis de escalada ou de disposição.}

%\clearpage
%\CGSetHeaderRight{VISÃO GERAL E HISTÓRIA}

%\CGDoubleTextSection{Visão Geral e História local}{
%A palavra "Cuó" é de origem indígena e significa acauã, ou seja, “pássaro devorador de cobras”. É o mesmo termo que deu origem ao nome do município de Caicó/RN, que significa \textit{Queicuó} e, traduzido, quer dizer “rio acauã”. No entanto, segundo o renomado historiador Luís da Câmara Cascudo (1968), Cuó significa divisão ou separação na língua nativa. A Serra do Cuó é separada por um boqueirão, dando-lhe um aspecto típico dividido, em uma região historicamente povoada pelos cariris.

%Hugolino D’Oliveira, recenseador no ano de 1940, já se referia à Serra do Cuó detalhando sua geografia: um pico com cerca de 428 metros de altura na ponta oriental da serra, formando o boqueirão que o separa do restante do maciço.
%}{}{
%Hoje a Serra do Cuó não abriga mais seus antigos habitantes. Uma nova geração se levantou e transformou a montanha em um formidável palco para os esportes de aventura. Suas trilhas, blocos e enormes paredões de pedra acolhem agora pessoas de vários lugares do Brasil e do mundo em busca de superação, liberdade e contato direto com a natureza.}

%\CGTextSingleColumn{A História da Escalada na Serra do Cuó}{
%As escaladas na Serra do Cuó tiveram seu início oficial em outubro de 2013, com a conquista das primeiras vias no Talhado dos Americanos. Brito Filho e Patrício Fernandes deram a primeira investida com a abertura de linhas esportivas. A via \textit{Boca Torta}, com proteção fixa, foi a pioneira, seguida imediatamente por outras duas fendas tradicionais com proteção em móvel e parada fixa.

%Após 2015, devido a alguns problemas de acesso ao setor dos blocões originais e às trilhas que partiam de lá, as atividades de desenvolvimento na serra sofreram um período de esfriamento. Mais recentemente, no entanto, o ímpeto de conquistas retornou com muita força, unindo novas gerações a antigos escaladores. Atualmente, o grande foco do montanhismo local está nos setores mais novos: Paraíso e Catedral.}

\clearpage
\CGOverlayImage{mapa_zonas_setores_left.jpg}
\CGOverlayImage{mapa_zonas_setores_right.jpg}
\clearpage
\CGSetHeaderRight{PLATAFORMAS DIGITAIS}

\CGTextSingleColumn{TheCrag}{
Este PDF é uma edição datada e não oficial. Consulte a comunidade no
\textbf{TheCrag} para atualizações das vias e contribua com correções,
comentários e registros de ascensões. As informações do PDF e do site não são
sincronizadas automaticamente; a data da edição não indica uma revisão completa
dos dados das vias.

O \textbf{TheCrag} é a principal plataforma global de escalada. É recomendado que você crie uma conta e registre suas ascensões (\textit{logbook}). Registrar seu progresso tem impactos reais:
\begin{itemize}
    \item Consolida a confirmação colaborativa do grau de dificuldade das vias.
    \item Permite comentários e descrições da via. Comunidade alerta sobre problemas nas proteções, marimbondos, ou quebras de agarras, etc.
    \item Contribuir para o registro da ``qualidade'' (estrelas) das vias, bem como fornece dados de quais vias têm mais repetições.
\end{itemize}
O uso do site torna a Serra do Cuó mais acessível para escaladores de fora. Escaneie o código abaixo (ou clique no link) para acessar a página oficial do pico:

\vspace{1em}
\CGQRCode[TheCrag: Serra do Cuó]{https://www.thecrag.com/en/climbing/brazil/area/14372948250}{serra_do_cuo_thecrag}
}
Como ajudar: cadastre as ascensões ou tentativas, comente descrições das vias e poste fotos das trilhas e vias.

\CGTextSingleColumn{Wikiloc}{
O \textbf{Wikiloc} é utilizado para mapear de forma precisa as trilhas de aproximação. Lá você encontrará as rotas de GPS, desde o acesso de carro até o estacionamento, além das trilhas específicas para cada zona e setor.

Utilize o código abaixo para acessar uma lista de trilhas da Serra do Cuó:

\vspace{1em}
\CGQRCode[Lista de Trilhas (Wikiloc)]{https://pt.wikiloc.com/wikiloc/user.do?id=20929616&event=favorites&listId=2256269}{serra_do_cuo_wikiloc}
}


% --- Zone Ingestion ---
\CGAddQRCodeToZoneHeader{https://www.thecrag.com/en/climbing/brazil/area/14373560118}{TheCrag: Zona Baixa}
\CGZoneHeader{zona_baixa_cover.jpg}{ZONA BAIXA}

\CGDoubleTextSection{}{
O muro de pedra é o setor de escalada mais próximo do estacionamento. Além desse bloco, ainda tem alguns blocos próximos, como o bloco do Machado, que tem grande potencial para exploração.
}{}{}

\clearpage
\CGSectorHeader{guide_purple}{Bloco Muro de Pedra}

\begin{CGSectorRoutes}\CGRoute{Muro de Berlim}{IV}{15m}{?+2}{Brito, Geraldo}
    \CGRoute{Alemanha Oriental}{IVsup}{\#m}{?+2}{Brito, Geraldo}[Parada compartilhada com Muro de Berlim]
    \CGRoute{Capitão Gondim}{VIsup}{\#m}{?+2}{Will}
    \CGRoute{Lista Negra}{VIsup}{\#m}{?+2}{Will}
    \CGRoute{A Cigana}{VIsup}{\#m}{?+2}{Will}
    \CGRoute{Tapuia}{VIsup}{\#m}{?+2}{Will}
    \CGRoute{Muralha da China}{VIsup}{\#m}{?+2}{Brito, Geraldo}
    \CGRoute{O Cigano}{VIIa}{\#m}{?+2}{Will, João}[Saida em movimento dinâmico, segue linha de buracos com agarras em ambos os lados.]
    \CGRoute{Falcão Acauã}{VIIb}{\#m}{?+2}{Will, João}
    \CGRoute{Jararaca}{VIIc}{\#m}{?+2}{Will}
    \CGRoute{Potiguar}{VIIc?}{\#m}{?+2}{Will}
    \CGRoute{Pedra de Responsa}{VI}{\#m}{?+2}{Brito, Geraldo, Bebeu}[Início em um V0 em regletes. A segunda metade não é de graça. ]
    \CGRoute{Pressão}{VIIa}{\#m}{?+2}{Brito}
    \CGRoute{Alemanha Ocidental}{VIIb}{\#m}{?+2}{Tomazini, Victor, Joel}[Inicio em um V3 delicado. Seguindo por agarras nos dois lados da aresta.]
    \CGRoute{Queda do Muro}{VIsup}{\#m}{?+2}{Tomazini, Victor, Joel}[Saída pela direita, mesmas chapas que a via `Alemanha Ocidental'.]
\end{CGSectorRoutes}


\CGOverlayImage{zona_baixa_map_muro.jpg}

\CGOverlayImage[H. O Cigano; I. Falcão Acauã; J. Jararaca; A. Muro de Berlim; B. Alemanha Oriental; C. Capitão Gondim; D. Lista Negra; E. A Cigana; L. Potiguar; M. Pedra de Responsa; N. Pressão; O. Alemanha Ocidental.  ]{zona_baixa_topo_muro.jpg}

\CGSectorHeader{guide_green}{Bloco do Machado}
\CGTextSingleColumn{}{Bloco com possibilidade de vias e boulders.}
\CGRouteImageRight[B. do Machado - Face Sul]{b_machado.jpg}{
    \begin{CGSectorRoutes}
    \end{CGSectorRoutes}
}
\clearpage

% --- paraiso.tex ---
\CGAddQRCodeToZoneHeader{https://www.thecrag.com/en/climbing/brazil/area/14373559746}{TheCrag: Paraíso}
\CGZoneHeader{paraiso_cover.jpg}{PARAÍSO}

\CGDoubleTextSection{}{Zona desenvolvida em 2026. Primeiras vias por Brito, Nayana, e Geraldo.  A primeira conquista do Paredão foi a via Mãe Natureza, e posteriormente a Águia no Ninho.
}{}{
A zona é dividida nos setores: Bloco da Macambira, Bloco do Portal, Bloco do Paraíso, Bloco da Fenda Linda, Bloco da Gameleira, Bloco da Gruta e Paredão Majestoso.
}%
\CGOverlayImage[Visão geral dos setores do Paraíso]{paraiso_map_sectors.jpg}
\clearpage

\CGSectorHeader{guide_red}{Bloco da Macambira}

O bloco mais ao Sul na trilha de acesso à zona do Paraíso.  Setor agradável com sombra o dia toda devido às arvores e ao bloco do Portal.%
\CGRouteImageRight[Bloco da Macambira]{paraiso_topo_macambira.png}{%
    \begin{CGSectorRoutes}%
        \noindent\CGRoute{Sem corda Ninguém escala}{IIIsup}{\#m}{}{conquistador(es)?}[Móvel]
        \CGRoute{Esqueci a parada}{IV}{\#m}{}{conquistador(es)?}[Mista]
    \end{CGSectorRoutes}
}
\clearpage

\CGSectorHeader{guide_purple}{Bloco do Portal}
Possui acesso principal passando pelo Portal do Paraíso e à esquerda do Bloco do Paraíso. Sendo este acesso desnivelado e com pedras soltas. Acesso alternativo pelo Bloco da Macambira, margeando o bloco do Portal pelo sudoeste. Ainda com potencial de expansão na face sudoeste e na aresta à direita da ``No Cuó dos outros é refresco''.

Bloco bastante arborizado e com sombra à tarde.%
\CGRouteImageRight[Brito na via `Treine, seu puto']{paraiso_topo_portal_3.png}{%
\begin{CGSectorRoutes}%
        \noindent\CGRoute{O Algoro}{VIIa}{\#m}{?+2}{conquistador(es)? (04/2026)}
        \CGRoute{Treine, seu puto}{VIsup}{\#m}{?+2}{conquistador(es)?}
        \CGRoute{Lucy in the sky}{V}{\#m}{?+2}{conquistador(es)? (04/2026)}
        \CGRoute{Cuó na mão}{VIIIa}{8m}{2+2}{Pedro, Marcel e Josias (07/2026)}
        \CGRoute{No Cuó dos outros é refresco}{VIsup}{\!\#m}{?+2}{Brito e Nayana (07/2026)}
    \end{CGSectorRoutes}
}
\CGFullWidthImage[Bloco do Portal]{paraiso_topo_portal.jpg}
\clearpage



\CGSectorHeader{guide_gray}{Bloco do Paraíso}
Bloco diretamente conectado com o do Portal. Permitindo escalar o dia todo na sombra alternando as faces.%
\CGRouteImageRight[Face Oeste - Ninho do Rouxinol e A Fúria]{paraiso_topo_bloco.png}{
    \begin{CGSectorRoutes}\CGRoute{Ninho do Rouxinol}{VI}{\#m}{?+2}{conquistador(es)?}[Face Oeste. ]
        \CGRoute{A Fúria}{VI}{\#m}{?+2}{Brito e Geraldo}[\CGstar1 Via com saida utilizando duas faces. Fúria até a segunda costura e um bom descanso. Segue em fenda de oposição.  ]
        \CGRoute{Seg também conquista}{Vsup}{\#m}{?+2}{Joel e Ingrid}[Face Norte]
        \CGRoute{Mole mesmo}{IIIsup}{\#m}{?+2}{conquistador(es)?}[Face Sul]
        \CGRoute{Arreboltamento}{Vsup}{\#m}{?+2}{Tomazini e Gilberto}[Face Sudeste]
    \end{CGSectorRoutes}
}
\clearpage

\CGSectorHeader{guide_cyan}{Bloco da Fenda Linda}

Bloco com sombra no ínicio da manhã, mas sem muita vegetação na face principal. Ao seu redor conta com algumas cavernas ``quase refrigeradas'' que permitem se abrigar do calor.  Ainda há possibilidade de expansão em outras faces desse bloco.

Conectado a esse bloco tem mais dois à Noroeste, estes ainda não explorados. Por meio do topo do bloco da Fenda Linda é possível acessar topo do bloco do meio.
\CGRouteImageRight[Via Fenda Linda]{paraiso_act_fenda.jpg}{
    \begin{CGSectorRoutes}\CGRoute{Fenda Linda}{Vsup}{\#m}{0+2}{conquistador(es)?}[Começa com uma saida atlética, e segue a fenda. Parede levemente positiva. Via em móvel com excelentes colocações. \\ \obs Acesso ao topo por escalaminhada à esquerda.]
        \CGRoute{Aprendendo a voar}{Vsup}{\#m}{?+2}{conquistador(es)?}
        \CGRoute{Urubu supersonico}{Vsup}{\#m}{?+2}{Brito}
        \CGRoute{Já era tarde}{III}{\#m}{3+2}{Breno Edon e Tomazini}
    \end{CGSectorRoutes}
}

\CGFullWidthImage[B. da Fenda Linda]{paraiso_topo_fenda.jpg}



\CGSectorHeader{guide_green}{Bloco da Gameleira}
Acesso principal ao bloco a partir do bloco da Fenda Linda, seguindo à direita até encontrar à Gameleira. Acesso passando pelos blocos à direita e contornando o bloco.

Acesso alternativo sem necessidade de passar por todos os setores do Paraíso.%
\CGRouteImageRight[Bloco da Gameleira]{paraiso_topo_gameleira.jpg}{%
    \begin{CGSectorRoutes}%
        \CGRoute{Purgatório}{VIIa}{12m}{3+2}{Gilberto e Tomazini (08/08/2026)}
        \CGRoute{Imunidade Solar}{VIIa}{15m}{3+2}{Gilberto e Tomazini (09/08/2026)}[\CGObs inacabada]
        \CGRoute{ODETE}{VI}{12m}{5+2}{Joel e Ingrid (09/08/2026)}
        \CGRoute{Tejucuóca}{V}{10m}{3+2}{Gilberto e Tomazini (08/08/2026)}
        \CGRoute{Cuó limpinho}{III}{6m}{1+2}{Gilberto e Tomazini (08/08/2026)}
        \CGRoute{Gameleira de Sangue}{V}{8m}{2+2}{Geraldo Filho e Yam Lima (26/08/2026)}[\CGObs confirmar ordem das vias.]
    \end{CGSectorRoutes}
}


\CGOverlayImage[Bloco da Gameleira]{paraiso_map_gameleira.jpg}
\CGOverlayImage[Blocos da Gameleira, Gruta e o Paredão Majestoso.]{paraiso_map_gruta.jpg}


\CGSectorHeader{guide_green}{Bloco da Gruta}
Bloco ainda sem vias abertas. O maior e mais alto bloco da zona do Paraíso. Tem três faces, uma face com negativo com abaulados, um face de rampão positivo e uma com uma gruta.%
\CGOverlayImage[Bloco da Gruta]{paraiso_cover_gruta.jpg}%
\CGSectorHeader{guide_orange}{Paredão Majestoso}%
Paredão imponente, com ganho de altitude de até 150 metros até a base e extensão de mais de 100m na sua parte mais alta. Parede característica vertical e de fendas horizontais. Ainda guarda grande potencial para exploração em vias de todas as dificuldades.

\begin{CGSectorRoutes}
    \CGRoute{Mãe Natureza}{6° VIIb E1}{70m}{14+2}{Brito e Geraldo}[3 enfiadas. A primeira via de varias enfiadas.]
    \CGRoute{O Ataque da Caranguejeira}{7° VIIb E1}{60m}{14+2}{Brito e Will}[Variante à esquerda da Águia no Ninho.]
    \CGRoute{Águia no Ninho}{7° VIIb E1}{60m}{14+2}{Brito e Will}[2 enfiadas. \obs Sem repetição.]
\end{CGSectorRoutes}

\CGFullWidthImage[Paredão Majestoso]{paraiso_topo_majestoso.jpg}





\CGZoneHeader{falesias_cover.jpg}{FALESIAS}
\CGDoubleTextSection{}{
Zona alta com parede contínua entre o paredão Majestoso do Paraiso e os pilares da Catedral. Possue acesso por trilha seguindo reto à Norte na bifurcação do bloco do Machado na trilha para a Catedral. Também possue acesso Norte, mas a trilha não está bem demarcada. Em periodos de chuva, existe um corrego e queda d'agua que cruza essa parede.
}{}{\CGObs Ainda não existe trilha demarcada conectando essa zona com a Catedral, existe um barranco devido ao corrego.

Zona ainda inexplorada. Zona sem nome.}


\CGAddQRCodeToZoneHeader{https://www.thecrag.com/en/climbing/brazil/area/14373559839}{TheCrag: Catedral}
\CGZoneHeader[ZONA DA]{catedral_cover.jpg}{CATEDRAL}


\CGDoubleTextSection{}{
Zona desenvolvida recentemente, iniciada no complexo dos ``blocões da Catedral'',  no setor da Bat Caverna. com as primeiras vias abertas durante a passagem de Will por Cuó no começo de 2026,  Bat Chaminé, Coringa e Robin.  Esse complexo teve rápido desenvolvimento durante o evento CuóRockTrip em Abril de 2026, com a maioria das vias dos setores da Bat Caverna, Corretor de pedra e Mansão Wayne foram conquistadas durante o evento.

}{}{Em Agosto, mais dois setores de escalada foram desenvolvidos na zona da Catedral. O setor das pinturas Rupestres teve inicio com a via  `La Dura', sem repetição, e posterior expansão margeando a rocha do complexo. Logo em seguida, em excursão com objetivo de conquistar a parede da Catedral, foi descoberto o setor do Vale Perdido, com acesso margeando a coluna direita da Catedral.}

\clearpage

\CGSectorHeader{guide_cyan}{Bat Caverna}

Setor com a maior concentração de vias de 4os e 5os. Acesso ao setor se dá pelo portal do inferno à esquerda das pinturas Rupestres.  Após passar pelo portal continua-se por margear o complexo de blocos até chegar no Corredor de Pedra. Seguindo mais um pouco chega à via Doce Ilusão, a mais à direita do setor.%
\begin{CGSectorRoutes}%
    \CGRoute{Bat Chaminé}{IVsup}{10m}{3+2}{Will}
    \CGRoute{Coringa}{V}{\#m}{6+2}{Will}
    \CGRoute{Robin}{V}{\#m}{4+2}{Will}
    \CGRoute{Pinguim}{Vsup}{\#m}{?+2}{Denn Malloy, Rytsa}[Mesma parada da Robin]
    \CGRoute{Charada}{Vsup}{\#m}{6+2}{Vini, Victor, Will}
    \CGRoute{BatMóvel}{IVsup}{\#m}{?+2}{Vini, Victor, Will}[Fenda]
    \CGRoute{Hera Venenosa}{V}{\#m}{4+2}{Tomazini, Izaú, Karla}
    \CGRoute{Madame Duarte}{IV}{\#m}{?+2}{Joel, Ingrid}
    \CGRoute{Madame Brito}{IV}{\#m}{?+2}{Jonas, Delrock}
    \CGRoute{Doce Ilusão}{IVsup}{\#m}{?+2}{Jonas, Delrock}
    \CGRoute{Gotham City}{VIIa}{\#m}{7+2}{Izaú, Will}
    \CGRoute{Rapadura-Quebra-Pedra}{Vsup}{\#m}{9+2}{Brito, Paulo, Geraldo}
    \CGRoute{Canto do Acauã}{IV}{30m}{3+2}{Brito, Paulo jr}[Mista]
\end{CGSectorRoutes}
\CGFullWidthImage[Bat Caverna, vista do Norte.]{catedral_topo_batcaverna.jpg}
\CGFullWidthImage[Parede direita da Bat Caverna.]{catedral_topo_batcaverna_direita.jpg}

\clearpage


\CGSectorHeader{guide_green}{{Face Leste - Blocões da Catedral}}
Parede Leste, acesso margeando a bat caverna pela esquerda, ou contornando por cima após o setor do Rupestre e Mansão Wayne. Setor ainda sem vias. Parede positiva com fendas. Permite acesso ao topo da Mansão Wayne e BatCaverna.%
\CGRouteImageRight[Face Leste - Blocões da Catedral]{catedral_cover_batcaverna.jpg}{
    \begin{CGSectorRoutes}
    \end{CGSectorRoutes}
}

\clearpage


\CGSectorHeader{guide_green}{Corredor das Pedras}
Setor característico de um corredor com pedregulhos enormes ao longo de sua extensão. Necessário fazer uma escalaminhada para acessar as vias `A Falta' e  `Serenidade', bem como a um mirante no topo com vista panorâmica.

Está direcionado ao norte e sem cobertura de árvore, fica no sol durante o inverno. Ainda assim, a parede esquerda (Corredor das 11) permanece na sombra no inicio da manhã, e a parede da direita está na sombra à tarde.

\begin{CGSectorRoutes}%
    \CGRoute{Corredor das 11}{VI}{\#m}{9+2}{Brito, Nayana}
    \CGRoute{A Flauta}{V}{\#m}{0+2}{Karla, Izaú}[Móvel]
    \CGRoute{Serenidade}{V}{\#m}{0+2}{Karla, Izaú}[Móvel]
    \CGRoute{Cuó é foda}{IVsup}{20m?}{0+2}{Aneke, Brito}[Móvel \#.3 ao \#3]
\end{CGSectorRoutes}

\CGOverlayImage[Vista de cima do complexo dos Blocões da Catedral. Setas indicam o Norte. Legenda: CP-Corredor de Pedra, R-Rupestres, MW-Mansão Wayne, L-Leste, BC-Bat Caverna.]{catedral_map_complexo.jpg}

\CGSectorHeader{guide_yellow}{Mansão Wayne}
Setor de acesso principal por dentro da Bat Caverna---o ``primeiro andar''. Essa passagem é por dentro da caverna. Acesso alternativo mais plano contornando após o setor Rupestres.

\begin{CGSectorRoutes}
%    \CGRoute{Batman 100 Perdão}{Proj}{\#m}{?+2}{Will, João}[Variante esquerda da Bruce Wayne]
    \CGRoute{Bruce Wayne}{Proj}{\#m}{?+2}{Pedro Caminha, Marcel Abdon}[8?]
    \CGRoute{Alfred, o mordomo}{Proj}{\#m}{6+2}{Will, João}
    \CGRoute{Comissário Gordon}{VIIa}{8m}{4+2}{Will, João}[Crux de V0/V1 na saída. Via negativa com diversos agarrões para mão e calcanhar. \\ \CGObs início na altura da primeira costura.]
    \CGRoute{Gota Serena}{VIIIa}{\#m}{?+2}{Josias, Denn Malloy, Buda}[Crux de V3 no negativo final.]
\end{CGSectorRoutes}
\CGOverlayImage[Mansão Wayne, via `Comissário Gordon']{catedral_topo_mansao_gordon.jpg}
\CGOverlayImage[Mansão Wayne, via `Gota Serena']{catedral_topo_mansao_gota.jpg}



\CGSectorHeader{guide_orange}{Rupestres}
\CGDoubleTextSection{}{
Setor com pinturas rupestres. Deu início em Agosto de 2026, com as vias `La Dura no Cuó' e `Aperta com Carinho' em uma visita de Aneke por Cuó.

Setor está na sombra durante toda a manhã, e possui base plana e com sombra de árvores e blocos possibilitando a escalada à tarde na sombra em algumas vias.
}{}{
O setor se conecta por trás do bloco da Mansão Wayne, cujo acesso é o mesmo que o topo das vias `Fura Olho' e `Equalização Perfeita'.
}%
\CGOverlayImage[Setor Rupestre, via `Aperta com carinho'.]{catedral_topo_rupestre_aperta.jpg}%
\begin{CGSectorRoutes}
    \CGRoute{Lado B Lado A}{VII?}{18m}{5+2}{Geraldo Filho e Yam Lima}[Mista. \textbf{Lado B:} Inicio tecnico, toda fixa, opção de descida pelo bloco. \textbf{Lado A:} segue em móvel. Jogo completo 0.1 a 4.]
    \CGRoute{La Dura no Cuó}{VIII?}{25m?}{8+2}{Aneke, Brito}[8b/c?]
    \CGRoute{Negativosa}{VIIb}{12m}{5+2}{Aneke, Brito}[Um lance levemente negativo técnico com regletes. Requer bom posicionamento de corpo. Segue em agarras generosas até o final.  ]
    \CGRoute{Aperta com Carinho}{Vsup}{12m}{4+2}{Aneke, Brito}[Atenção na segunda costurada. Segue com varias agarras e muitas opções de descanso. ]
    \CGRoute{Caminho Rupestre}{V}{\#m}{0+2}{Aneke, Brito}[Móvel .3 ao \#4]
    \CGRoute{Variante Diedro}{V}{\#m}{0+2}{Aneke, Brito}[Mesma parada que a ``Caminho Rupestre''. Segue na fenda do diedro à direita. ]
    \CGRoute{Movelandia}{VI}{\#m}{1+2}{Aneke, Thibor}[\#.3 ao \#4]
    \CGRoute{Árvore morta}{VI?}{\#m}{1+2}{Aneke, Thibor}[\#.3 ao \#4. \CGObs sem repetição. ]
    \CGRoute{Fura Olho}{VI}{\#m}{1}{Vini, Lu, Thibor}[\CGObs Proteções Pendentes (2+2). Saída em um batente raso. Rocha com fendas diagonais, segue com opções de agarras à esquerda.]
    \CGRoute{Equalização Perfeita}{V}{\#m}{0}{Vini, Lu, Thibor}[\CGObs  Proteções Pendentes (2+2). Saída na fenda aberta e segue reto e levemente para esquerda com agarras suficientes. ]
\end{CGSectorRoutes}

\CGOverlayImage[Vista de entrada do setor Rupestre.]{catedral_cover_rupestre.jpeg}
\clearpage


\CGSectorHeader{guide_red}{Vale Perdido}
Primeiro setor desenvolvido na Catedral. O acesso para o Vale é pelo pilar direito, localizado à Oeste. O setor permite escalada em todas as faces, e as conquistas estão só começando.

\begin{CGSectorRoutes}
    \CGRoute{Cuózinho}{3 IVsup}{50m}{?+2}{Aneke, Jonas}[\#4, \#.75]
    \CGRoute{Eu chego lá}{VIII?}{15m}{5+2}{Brito Filho e Yam Lima (22/08/2026)}[\CGObs sem repetições]
    \CGRoute{Murphy}{VIIb}{\#m}{?+2}{Aneke, Del}[\CGstar1 via recomendada]
    \CGRoute{Som da Laca}{VI}{25m?}{6+2}{conquistador(es)?}
    \CGRoute{São Geraldo}{IV}{20m?}{5+2}{Aneke, Jonas}
    \CGRoute{Legalize It}{Vsup}{25m?}{7+2}{DelRock, Jonas, Aneke}[\CGstar1 via recomendada]
    \CGRoute{Cuó Solitário}{VI}{\#m}{3+2}{Aneke}
\end{CGSectorRoutes}



\CGOverlayImage[Catedral - Vale Perdido.]{catedral_map_vale.jpg}


\clearpage

\CGAddQRCodeToZoneHeader{https://www.thecrag.com/en/climbing/brazil/area/14373559644}{TheCrag: Talhado dos Americanos}
\CGZoneHeader{talhado_act_boca_torta.jpg}{TALHADO DOS AMERICANOS}
\CGDoubleTextSection{}{
    Setor com as primeiras vias do Cuó.
}{}{}
\clearpage

\CGSectorHeader{guide_cyan}{Setor da Entrada}
\CGRouteImageRight[Conquista Boca Torta.]{talhado_act_boca_torta.jpg}{
    \begin{CGSectorRoutes}
        \CGRoute{Free Solo}{IIsup}{\#m}{?+2}{conquistador(es)?}[Free solo]
        \CGRoute{Via do Bloco}{Proj}{\#m}{?+2}{conquistador(es)?}[Nome/grau não confirmados]
        \CGRoute{Boca Torta}{VI}{\#m}{?+2}{Brito Filho e Patrício Fernandes}[Primeira via esportiva do Cuó]
        \CGRoute{Fenda das Irmãs}{V}{\#m}{Móvel}{conquistador(es)?}
        \CGRoute{Fendona}{V}{\#m}{Móvel}{conquistador(es)?}
    \end{CGSectorRoutes}
}
\clearpage

\CGSectorHeader{guide_green}{Parte Alta}
\CGTextSingleColumn{}{Fica na face leste, caracterizada por trechos negativos e vias mais exigentes.}
\begin{CGSectorRoutes}
    \CGRoute{Inacabada}{Proj}{\#m}{3 chapas}{conquistador(es)?}[Via inacabada]
    \CGRoute{Du CuóRalho}{VIIa}{\#m}{?+2}{conquistador(es)?}
    \CGRoute{O Ronco}{VIsup}{\#m}{?+2}{conquistador(es)?}
    \CGRoute{Mr. Ozzy}{V}{\#m}{?+2}{conquistador(es)?}
    \CGRoute{CuóRalina}{Vsup}{\#m}{?+2}{conquistador(es)?}
    \CGRoute{Bicho Pau}{VIsup}{\#m}{?+2}{conquistador(es)?}
    \CGRoute{Alucicrazy}{VIIb}{\#m}{?+2}{conquistador(es)?}
    \CGRoute{Nona Sinfonia}{VIII?}{\#m}{?+2}{conquistador(es)?}
    \CGRoute{Aryado}{VIII?}{\#m}{?+2}{conquistador(es)?}
\end{CGSectorRoutes}
\clearpage

\CGSectorHeader{guide_yellow}{Setor Árvore da Vida}
\begin{CGSectorRoutes}
    \CGRoute{Ladrão de F.A.}{VI}{\#m}{?+2}{conquistador(es)?}
    \CGRoute{Árvore da vida}{IVsup}{\#m}{?+2}{conquistador(es)?}
    \CGRoute{Todos por um}{Vsup}{\#m}{?+2}{conquistador(es)?}
    \CGRoute{Faltou Cuó}{V}{\#m}{?+2}{conquistador(es)?}
    \CGRoute{São Longuinho é foda}{V}{\#m}{?+2}{conquistador(es)?}
\end{CGSectorRoutes}
\clearpage

\CGSectorHeader{guide_orange}{Ceisão}
\begin{CGSectorRoutes}
    \CGRoute{Ceisão}{VIsup}{\#m}{?+2}{Brito e Aneke}
    \CGRoute{Olho gordo}{V}{\#m}{?+2}{Brito e Aneke}
    \CGRoute{Jacaronte}{V}{\#m}{?+2}{Brito, Patrício e Aneke}
    \CGRoute{Chaminé blackout}{IIsup}{\#m}{?+2}{Brito e Patrício}
\end{CGSectorRoutes}
\clearpage


% --- Deferred Index Rendering ---

% Endmatter: Geology, History, and Indexes

\CGBackmatterChapter[2]{História da Conquista}{
    ...
    %A exploração desta zona começou com a ascensão de vias clássicas na década de 90.
    %A ética de conquista local prioriza o uso de proteções móveis onde a rocha permite.
}

\CGBackmatterChapter[1]{Geologia}{
    ...
}

\clearpage
\CGSetHeaderRight{Índice Alfabético}
\phantomsection
\addcontentsline{toc}{section}{Índice Alfabético}
\CGTextSingleColumn{Índice Alfabético}{}
\CGPrintAlphaList

\clearpage
\CGSetHeaderRight{Índice por Dificuldade}
\phantomsection
\addcontentsline{toc}{section}{Índice por Dificuldade}
\CGTextSingleColumn{Índice por Dificuldade}{}
\CGPrintGradeList

\clearpage
\CGSetHeaderRight{Trilhas e Acesso}
\phantomsection
\addcontentsline{toc}{section}{Apêndice: Trilhas e Acesso (Wikiloc)}

\vspace*{2em}\noindent
{\color{brand_color}\CGSectorTitleFont Apêndice: Trilhas e Acesso}\par\vspace{1.5em}

{\CGBodyFont 
As trilhas de aproximação para os setores da Serra do Cuó podem ser confusas para visitantes de primeira viagem devido aos blocos e vegetação densa. Abaixo estão os códigos de acesso direto para as rotas de GPS no \textbf{Wikiloc}.
}
\vspace{2em}

\begin{multicols}{2}
    \textbf{Acesso Geral (Estacionamento)}\par\vspace{4pt}
    \CGQRCode[Rota de Carro]{https://pt.wikiloc.com/wikiloc/user.do?id=20929616&event=favorites&listId=2256269}\par
    \vspace{2em}
    
    \textbf{Zona Baixa \& Catedral}\par\vspace{4pt}
    \CGQRCode[Trilha Principal]{https://pt.wikiloc.com/wikiloc/user.do?id=20929616&event=favorites&listId=2256269}\par
    \vspace{2em}

    \columnbreak

    \textbf{Paraíso \& Setores Altos}\par\vspace{4pt}
    \CGQRCode[Boqueirão e Paraíso]{https://pt.wikiloc.com/wikiloc/user.do?id=20929616&event=favorites&listId=2256269}\par
    \vspace{2em}
\end{multicols}



\end{document}
